

On
peut donc de façon évidente étendre @@M0@@ en un homomorphisme de @@M1@@-algèbres : @@M2@@. On a donc :
@@M3@@

\textbf{Définition 6.4.} \uline{L'homomorphisme} @@M4@@ \uline{qui vient d'être} \uline{défini s'appelle caractère de Chern}.

Si @@M5@@ est de caractéristique @@M6@@ (i.e. est une @@M7@@-algèbre), le caractère de Chern est un isomorphisme~: en effet, soit @@M8@@~; on peut écrire @@M9@@, où @@M10@@ est sans terme de degré nul~; on définit alors un morphisme inverse de @@M11@@ par @@M12@@.

\textbf{Proposition 6.5.} \uline{Soit} @@M13@@ \uline{une} @@M14@@-\uline{algèbre graduée}~; \uline{considérons sur} @@M15@@ \uline{la filtration définie comme suit} (\uline{filtration par le degré}) :
@@M16@@
\uline{où} @@M17@@. \uline{On a les congruences suivantes} :

\begin{enumerate}[label=\roman*)]
\item \uline{Soient} @@M18@@, \uline{et} @@M19@@~; \uline{posons} @@M20@@~; \uline{alors} :
@@M21@@

\item \uline{Soient} @@M22@@, @@M23@@~; \uline{posons}
@@M24@@, @@M25@@~; \uline{alors} :
@@M26@@
\uline{et par suite}, \uline{la filtration par le degré est une filtration d'anneau}.
\end{enumerate}

i) Le caractère de Chern est par construction un morphisme de foncteurs filtrés~; de la congruence @@M27@@, on déduit
